\chapter{Event Schema}
\label{app:schema}

This appendix lists everything the editor records. The source of truth is
\nolinkurl{backend/studentide/schema.py}, which the server uses to validate every
uploaded event; \cref{tab:event-schema} is generated from it by
\nolinkurl{backend/scripts/build_schema_appendix.py}, and a test fails if the two
drift apart. A prose version for developers is \nolinkurl{docs/event-schema.md}.

\section{Envelope and storage}
A session's events are uploaded in batches that name their schema version. Each
event has the same envelope:
\begin{description}[nosep,font=\ttfamily\mdseries]
  \item[seq] a per-session counter from 0, assigned in the browser. Uploads are
  retried until acknowledged and the server ignores a repeated \texttt{seq}, so
  no event is stored twice; a gap means a lost event.
  \item[t] milliseconds since the session started, from the browser's monotonic
  event clock. Timing features use this, never server time. Sort by
  \texttt{seq} for causal order and by \texttt{t} for timing.
  \item[type] one of the types in \cref{tab:event-schema}.
  \item[data] the payload for that type, validated against its model; unknown
  fields are refused.
\end{description}

Events are stored in SQLite in an \texttt{events} table keyed by
\texttt{(session\_id, seq)}. A \texttt{sessions} row holds the student, the
lab and question (with a copy of the question as the student saw it,
including its language), the sitting identifier, the context that marks non-study sessions
(\texttt{learn}, \texttt{preview}; \cref{sec:non-study}), client metadata
(browser or desktop app version, platform, screen) and, once raw keys are
deleted, the time they were purged. A \texttt{runs} table holds each run's full
code, input, output and error type, and \texttt{session\_windows} holds the
window summaries that replace raw keys (\cref{sec:privacy}).

\section{Versions}
Every version is additive: a batch valid under an older version is valid under
every newer one, so recordings from the first pilot day are still read by the
current code.
\begin{description}[nosep]
  \item[v1] the browser editor: \texttt{key}, \texttt{edit},
  \texttt{clipboard}, \texttt{cursor}, \texttt{run}, \texttt{snapshot},
  \texttt{focus}, \texttt{session\_start}, \texttt{session\_end}.
  \item[v2] the desktop app's Sandbox Labs: where a key was typed (editor,
  terminal, an editor inside the terminal), the file of a multi-file workspace,
  runs typed in the terminal, and terminal commands, output, pastes, package
  installs and the hosts a lab connected to. Only editor keys feed the
  detector.
  \item[v3] files changed outside the editor and ports opened in a Sandbox Lab;
  the desktop app also signs its batches and chains them by hash, so a
  rewritten or missing batch shows as a record break in the review.
  \item[v4] the paste guard (\cref{sec:paste-guard}): a paste or drop records
  where its text came from (\texttt{editor}, \texttt{page}, \texttt{document} or
  \texttt{external}) and whether it was blocked; allowed in-IDE pastes are edits
  with origin \texttt{input.paste.internal}.
\end{description}
None of the v2 and v3 Sandbox Lab events is read by the detector evaluated in
this thesis; they feed the review's evidence lanes only.

\section{Payloads}
% Generated by backend/scripts/build_schema_appendix.py from backend/studentide/schema.py.
% Do not edit by hand; rerun the script after a schema change.
\begingroup\small
\begin{longtable}{@{}p{0.2\linewidth}p{0.5\linewidth}p{0.22\linewidth}@{}}
\caption{Event payloads by type (schema version 4). Every event also carries \texttt{seq}, \texttt{t} and \texttt{type}.}\label{tab:event-schema}\\
\toprule Field & Type & Default \\ \midrule \endfirsthead
\toprule Field & Type & Default \\ \midrule \endhead
\bottomrule \endfoot
\multicolumn{3}{@{}p{0.95\linewidth}}{\texttt{key}: A physical key press or release. Timing source for keystroke dynamics.} \\*
\quad\texttt{action} & \texttt{"down" | "up"} & \texttt{required} \\
\quad\texttt{code} & \texttt{str} & \texttt{required} \\
\quad\texttt{cls} & \texttt{"letter" | "digit" | "space" | "enter" | "tab" | "punct" | "backspace" | "delete" | "nav" | "modifier" | "other"} & \texttt{required} \\
\quad\texttt{repeat} & \texttt{bool} & \texttt{false} \\
\quad\texttt{shift} & \texttt{bool} & \texttt{false} \\
\quad\texttt{ctrl} & \texttt{bool} & \texttt{false} \\
\quad\texttt{alt} & \texttt{bool} & \texttt{false} \\
\quad\texttt{meta} & \texttt{bool} & \texttt{false} \\
\quad\texttt{surface} & \texttt{"editor" | "terminal" | "terminal-editor"} & \texttt{"editor"} \\
\quad\texttt{path} & \texttt{str | null} & \texttt{null} \\
\midrule
\multicolumn{3}{@{}p{0.95\linewidth}}{\texttt{edit}: One editor transaction that changed the document.} \\*
\quad\texttt{origin} & \texttt{str} & \texttt{required} \\
\quad\texttt{changes} & \texttt{list[tuple[int, int, str]]} & \texttt{required} \\
\quad\texttt{len\_before} & \texttt{int} & \texttt{required} \\
\quad\texttt{len\_after} & \texttt{int} & \texttt{required} \\
\quad\texttt{path} & \texttt{str | null} & \texttt{null} \\
\midrule
\multicolumn{3}{@{}p{0.95\linewidth}}{\texttt{clipboard}: paste / copy / cut / drop on the editor. The pasted text itself arrives in the following edit event (origin input.paste); this records the clipboard side.} \\*
\quad\texttt{action} & \texttt{"paste" | "copy" | "cut" | "drop"} & \texttt{required} \\
\quad\texttt{length} & \texttt{int} & \texttt{required} \\
\quad\texttt{lines} & \texttt{int} & \texttt{required} \\
\quad\texttt{text} & \texttt{str} & \texttt{required} \\
\quad\texttt{truncated} & \texttt{bool} & \texttt{false} \\
\quad\texttt{source} & \texttt{"editor" | "page" | "document" | "external" | null} & \texttt{null} \\
\quad\texttt{blocked} & \texttt{bool} & \texttt{false} \\
\quad\texttt{surface} & \texttt{"editor" | "terminal" | "terminal-editor"} & \texttt{"editor"} \\
\midrule
\multicolumn{3}{@{}p{0.95\linewidth}}{\texttt{cursor}: Selection change not caused by an edit (clicks, arrow keys). Throttled.} \\*
\quad\texttt{anchor} & \texttt{int} & \texttt{required} \\
\quad\texttt{head} & \texttt{int} & \texttt{required} \\
\quad\texttt{line} & \texttt{int} & \texttt{required} \\
\midrule
\multicolumn{3}{@{}p{0.95\linewidth}}{\texttt{run}: The student pressed Run. Full result lives in the runs table under run\_id.} \\*
\quad\texttt{run\_id} & \texttt{str} & \texttt{required} \\
\quad\texttt{code\_sha256} & \texttt{str} & \texttt{required} \\
\quad\texttt{exit\_code} & \texttt{int | null} & \texttt{null} \\
\quad\texttt{duration\_ms} & \texttt{float | null} & \texttt{null} \\
\quad\texttt{error\_type} & \texttt{str | null} & \texttt{null} \\
\quad\texttt{source} & \texttt{"button" | "test" | "terminal" | null} & \texttt{null} \\
\quad\texttt{command} & \texttt{str | null} & \texttt{null} \\
\midrule
\multicolumn{3}{@{}p{0.95\linewidth}}{\texttt{snapshot}: Full document text, taken at session start, every N seconds, on run and at end. Lets analysis jump into the middle of a session and double-checks replay.} \\*
\quad\texttt{reason} & \texttt{"start" | "periodic" | "run" | "end"} & \texttt{required} \\
\quad\texttt{text} & \texttt{str} & \texttt{required} \\
\quad\texttt{path} & \texttt{str | null} & \texttt{null} \\
\midrule
\multicolumn{3}{@{}p{0.95\linewidth}}{\texttt{focus}: Editor focus, window focus, and tab visibility changes. In a lab with several questions, target "question" marks the student switching to another question of the same sitting ("out") and back to this one ("in"), and in an any-language competition the student switching this question to another language (each language is its own session) and back.} \\*
\quad\texttt{target} & \texttt{"editor" | "window" | "page" | "question"} & \texttt{required} \\
\quad\texttt{state} & \texttt{"in" | "out"} & \texttt{required} \\
\midrule
\multicolumn{3}{@{}p{0.95\linewidth}}{\texttt{session\_start}, \texttt{session\_end}: session\_start / session\_end markers.} \\*
\quad\texttt{wall\_clock\_ms} & \texttt{float} & \texttt{required} \\
\quad\texttt{reason} & \texttt{str | null} & \texttt{null} \\
\midrule
\multicolumn{3}{@{}p{0.95\linewidth}}{\texttt{term\_cmd} (since v2): v2: a command line the desktop terminal's shell ran, from its shell hooks (best effort: the student has sudo and can remove them). phase "start" when it begins, "end" when the prompt is back, with the exit code and duration.} \\*
\quad\texttt{term\_id} & \texttt{str} & \texttt{required} \\
\quad\texttt{phase} & \texttt{"start" | "end"} & \texttt{required} \\
\quad\texttt{command} & \texttt{str} & \texttt{required} \\
\quad\texttt{cwd} & \texttt{str} & \texttt{required} \\
\quad\texttt{truncated} & \texttt{bool} & \texttt{false} \\
\quad\texttt{exit\_code} & \texttt{int | null} & \texttt{null} \\
\quad\texttt{duration\_ms} & \texttt{float | null} & \texttt{null} \\
\midrule
\multicolumn{3}{@{}p{0.95\linewidth}}{\texttt{term\_output} (since v2): v2: output of one terminal command: bytes written and a clipped, escape-free tail.} \\*
\quad\texttt{term\_id} & \texttt{str} & \texttt{required} \\
\quad\texttt{bytes} & \texttt{int} & \texttt{required} \\
\quad\texttt{tail} & \texttt{str} & \texttt{required} \\
\quad\texttt{truncated} & \texttt{bool} & \texttt{false} \\
\midrule
\multicolumn{3}{@{}p{0.95\linewidth}}{\texttt{term\_paste} (since v2): v2: a paste into the terminal, like ClipboardData for the editor.} \\*
\quad\texttt{term\_id} & \texttt{str} & \texttt{required} \\
\quad\texttt{length} & \texttt{int} & \texttt{required} \\
\quad\texttt{lines} & \texttt{int} & \texttt{required} \\
\quad\texttt{text} & \texttt{str} & \texttt{required} \\
\quad\texttt{truncated} & \texttt{bool} & \texttt{false} \\
\quad\texttt{sha256} & \texttt{str} & \texttt{required} \\
\midrule
\multicolumn{3}{@{}p{0.95\linewidth}}{\texttt{pkg\_install} (since v2): v2: a package command that succeeded in the student's instance (feeds the setup recipe).} \\*
\quad\texttt{manager} & \texttt{str} & \texttt{required} \\
\quad\texttt{action} & \texttt{"install" | "remove"} & \texttt{required} \\
\quad\texttt{packages} & \texttt{list[str]} & \texttt{required} \\
\quad\texttt{command} & \texttt{str} & \texttt{required} \\
\midrule
\multicolumn{3}{@{}p{0.95\linewidth}}{\texttt{net\_conn} (since v2): v2: the instance reached a new host through the lab's proxy. Host name only, never contents.} \\*
\quad\texttt{host} & \texttt{str} & \texttt{required} \\
\quad\texttt{port} & \texttt{int | null} & \texttt{null} \\
\quad\texttt{allowed} & \texttt{bool} & \texttt{required} \\
\midrule
\multicolumn{3}{@{}p{0.95\linewidth}}{\texttt{fs\_change} (since v3): v3: a workspace file changed and our editor did not write it, seen by the sidecar watcher (root in the instance cannot stop it). origin: "terminal-editor" (vim/nano was in the foreground), "terminal" (another command was running: cat >, git checkout, a script), "outside" (no command we can see: docker cp, a host tool, a background job). added = the new lines of this save (clipped), the evidence check reads it.} \\*
\quad\texttt{path} & \texttt{str} & \texttt{required} \\
\quad\texttt{size} & \texttt{int} & \texttt{required} \\
\quad\texttt{sha256} & \texttt{str} & \texttt{required} \\
\quad\texttt{origin} & \texttt{"terminal-editor" | "terminal" | "outside"} & \texttt{required} \\
\quad\texttt{command} & \texttt{str | null} & \texttt{null} \\
\quad\texttt{deleted} & \texttt{bool} & \texttt{false} \\
\quad\texttt{binary} & \texttt{bool} & \texttt{false} \\
\quad\texttt{added} & \texttt{str} & \texttt{""} \\
\quad\texttt{added\_lines} & \texttt{int} & \texttt{0} \\
\quad\texttt{removed\_lines} & \texttt{int} & \texttt{0} \\
\quad\texttt{diff} & \texttt{str} & \texttt{""} \\
\quad\texttt{truncated} & \texttt{bool} & \texttt{false} \\
\midrule
\multicolumn{3}{@{}p{0.95\linewidth}}{\texttt{port\_open} (since v3): v3: a program in the instance started (open) or stopped listening on a TCP port.} \\*
\quad\texttt{port} & \texttt{int} & \texttt{required} \\
\quad\texttt{process} & \texttt{str} & \texttt{""} \\
\quad\texttt{open} & \texttt{bool} & \texttt{required} \\

\end{longtable}
\endgroup


\paragraph{Not an event.} The staff-only note of how a conversation went
(\cref{sec:review-timeline}) is stored in its own table, keyed by session,
channel and the moment's estimated start. It is never part of a session's event
stream, never read by the detector, and no student endpoint returns it.
